\subsection{PARTIAL SUMS; ASSOCIATIVITY}

PROPOSITION 2. In a complete group G, every subfamily of a summable family is summable.

For if Cauchy's criterion is satisfied by a family \((x_{i})_{i\in I}\) , then it is trivially satisfied by every subfamily.

If \((x_{i})_{i\in I}\) is summable, it follows therefore that the sum \(\sum_{i\in J}x_{i}\) is defined for every subset \(J\) of \(I\) , finite or not: it is again called the partial sum of the family \((x_{i})\) , corresponding to the subset \(J\) of the index set. The set of partial sums of a summable family is evidently contained in the closure of the set of finite partial sums.

THEOREM 2 (Associativity of the sum). Let \((x_{i})_{i\in I}\) be a summable family in a complete group \(\mathbf{G}\) , and let \((I_{\lambda})_{\lambda \in L}\) be any partition of \(\mathbf{I}\) . If \(s_{\lambda}\) denotes \(\sum_{i\in I_{\lambda}}x_{i}\) , then the family \((s_{\lambda})_{\lambda \in L}\) is summable and has the same sum as the family \((x_{i})_{i\in I}\) .

Let \(s = \sum_{i \in I} x_i\) , and let \(V\) be any closed neighbourhood of 0 in \(G\) . Then there is a finite subset \(J_0\) of \(I\) such that, for each finite subset \(J\) of \(I\) which contains \(J_0\) , we have \(\sum_{i \in J} x_i \in s + V\) . Let \(K_0\) be the subset of \(L\) consisting of indices \(\lambda\) such that \(J_\lambda = I_\lambda \cap J_0\) is not empty: \(K_0\) is clearly finite. Let \(K\) be any finite subset of \(L\) which contains \(K_0\) ; we shall show that \(\sum_{\lambda \in K} s_\lambda \in s + V\) , which will establish the theorem. Now \(s_\lambda\) is very close to a finite partial sum of \((x_i)\) , whose indices all belong to \(I_\lambda\) ; more precisely, given any symmetric neighbourhood \(W\) of 0, there exists for each \(\lambda \in K\) a finite subset \(H_\lambda\) of \(I_\lambda\) , containing \(J_\lambda\) and such that \(s_\lambda - \sum_{i \in H_\lambda} x_i \in W\) . Let \(J = \bigcup_{\lambda \in K} H_\lambda\) ; then \(J\) is a finite subset of \(I\) containing \(J_0\) , and we have

\[
\sum_ {i \in J} x _ {i} = \sum_ {\substack {i \in \bigcup_ {\lambda \in K} H _ {\lambda}}} x _ {i} = \sum_ {\lambda \in K} (\sum_ {i \in H _ {\lambda}} x _ {i})
\]

by the associativity of finite sums. Hence by reason of the choice of \(J_{0}\) and the \(H_{\lambda}\) , we have

\[
\sum_ {\lambda \in \mathbf {K}} s _ {\lambda} \in s + \mathbf {V} + n \mathbf {W}
\]

where \(n\) is the number of elements in \(\mathbf{K}\) ; this relation holds for each \(\mathbf{W}\) , hence also \(\sum_{\lambda \in \mathbb{K}} s_{\lambda} \in s + \mathbf{V}\) , since \(\mathbf{V}\) (being closed) is the intersection of the neighbourhoods \(\mathbf{V} + n\mathbf{W}\) {[}§ 3, no. 1, formula (1){]}.

We can thus write the associativity formula for sums:

\[
\sum_ {\lambda \in L} (\sum_ {i \in I _ {\lambda}} x _ {i}) = \sum_ {i \in \bigcup_ {\lambda \in L} I _ {\lambda}} x _ {i},\tag{1}
\]

which is valid whenever the family \((I_{\lambda})\) is a partition of its union and the right-hand side is defined. In particular, if the index set is a product \(I = L \times M\) , and if the ``double'' family \((x_{\lambda\mu})_{(\lambda,\mu) \in L \times M}\) is summable, we have the formula of change of order of summation

\[
\sum_ {(\lambda , \mu) \in \mathbf {L} \times \mathbf {M}} x _ {\lambda \mu} = \sum_ {\lambda \in \mathbf {L}} (\sum_ {\mu \in \mathbf {M}} x _ {\lambda \mu}) = \sum_ {\mu \in \mathbf {M}} (\sum_ {\lambda \in \mathbf {L}} x _ {\lambda \mu}).\tag{2}
\]

It should be remarked that the left-hand side of (1) can have a meaning, without the right-hand side being defined. Consider for example the case where \(I = L \times \{1, 2\}\) and L is infinite, and \(I_{\lambda}\) consists of the two elements \((\lambda, 1)\) and \((\lambda, 2)\) ; if we take \(x_{\lambda,1} = a\) , \(x_{\lambda,2} = -a\) , where \(a\) is any non-zero element of \(\mathbf{G}\) , then all the partial sums corresponding to the \(I_{\lambda}\) are zero, and therefore the left-hand side of (1) is defined and equal to 0, whereas the right-hand side of (1) has no meaning, as Proposition 1 shows.

Likewise, it can happen that the left-hand side of (2) is not defined but that each of the two ``double sums'' in (2) has a meaning; and the elements of G which they represent need not be equal (see Chapter IV, § 7, Exercise 17).

Thus although it is always possible to ``associate'' the terms of a sum, it is not possible, on the other hand, to ``dissociate'' into their elements those terms of a sum which appear themselves as sums. Nevertheless, this operation is legitimate whenever the number of these ``dissociable'' terms is finite.

PROPOSITION 3. Let \((x_{\iota})_{\iota \in \mathbf{I}}\) be a family of points of a group \(\mathbf{G}\) , and let \((I_{\lambda})_{\lambda \in \mathbf{L}}\) be a finite partition of \(\mathbf{I}\) . If each of the subfamilies \((x_{\iota})_{\iota \in \mathbf{I}_{\lambda}}\) is summable, then the family \((x_{\iota})_{\iota \in \mathbf{I}}\) is summable and the formula (1) is valid.

It is enough to prove the proposition when \(L = (1,2)\) ; having done so, we then proceed by induction on the number of elements of L. Let \(s_{1} = \sum_{i \in I_{1}} x_{i}\) and \(s_{2} = \sum_{i \in I_{2}} x_{i}\) . For each neighbourhood V of the origin, there is a finite subset \(J_{1}\) (resp. \(J_{2}\) ) of \(I_{1}\) (resp. \(I_{2}\) ) such that, for each finite subset \(H_{1}\) (resp. \(H_{2}\) ) of \(I_{1}\) (resp. \(I_{2}\) ) containing \(J_{1}\) (resp. \(J_{2}\) ), we have \(\sum_{i \in H_{1}} x_{i} \in s_{1} + V\) (resp. \(\sum_{i \in H_{2}} x_{i} \in s_{2} + V\) ). If we put \(J_{0} = J_{1} \cup J_{2}\) , it follows that, for each finite subset H of I which contains \(J_{0}\) , we have \(\sum_{i \in H} x_{i} \in s_{1} + s_{2} + V + V\) , and the result follows.

